\section{Framed formal monodromy: a conditional refinement}
\label{sec:framed-monodromy}

Section~\ref{sec:atomic-one-step} proved the one-stabilization theorem using
one ordinary Hodge atom.  We now ask for a stronger statement about the
entire numerical small even quantum connection: whether its
primitive-sixth framed formal-monodromy multiplicity satisfies direct
operation formulas under projective bundles and blowups.  The actual
Iritani and Iritani--Koto comparison isomorphisms preserve the original
loop framing after coefficient-field extension.  The string equation and
divisor equation also remove, exactly, the scalar degree-zero term and any
fixed degree-two term whose Novikov character descends to the coefficient
ring.  What remains is narrower: a residual reconstruction-tail displacement,
and a logically separate specialization step
in the divisor-tagging argument.  We isolate these as Hypotheses~5.7R
and~5.7T below.

The results of this section are therefore logically separate from
Theorem~\ref{thm:every-cubic}.  Definition~\ref{def:framed-sixth-multiplicity},
the cubic computation, and the product-with-projective-space formula below
are unconditional.  The projective-bundle and blowup operation formulas
use Hypothesis~5.7R.  We prove specialized low-dimensional vanishing
directly, without Hypothesis~5.7T, for nef-canonical targets, for
\(\PP^1\) and \(\PP^2\), and---assuming only Hypothesis~5.7R---for ruled
surfaces over curves of positive genus.  In the present proof
Hypothesis~5.7T is used only for the remaining rational ruled and nonminimal
surface cases.  Consequently the resulting birational invariance of
\(\nu_6\), and the identification \(\nu_6(V)=2\) for genus-eight Fano
threefolds, remain conditional on both hypotheses.


Proposition~\ref{prop:residue-discriminant-exponents} explains how the two
quantum invariants in this paper are related on the rank-two block that
matters for the cubic.  The atomic invariant $\delta^\sharp$ is the square of
the difference of the residue representatives selected by the canonical
elementary modification, and therefore forgets a common scalar shift.  The
framed invariant $\nu_6$ is finer: it tests the absolute cyclotomic placement
of the formal exponent classes, together with the original-loop descent
data when ramification is present.  For the cubic zero block, Section~\ref{sec:atomic-one-step}
already finds the representatives $-1/6$ and $-5/6$; the calculation below
recovers those same exponents directly from formal solutions and records
that their classes give the two primitive sixth roots.  Thus Sections~4 and
5 use two levels of the same formal data rather than unrelated numerical
obstructions.

Let \(T\) be a smooth projective variety.  Write
\[
 \Lambda_T=\C[[\operatorname{NE}^{\mathbf Z}_{\rm num}(T)]]
\]
for the completed monoid ring of the image of effective homology classes in
the numerical curve lattice
\[
 N_1(T)=\operatorname{im}\!\left(
 H_2(T,\mathbf Z)\longrightarrow
 \operatorname{Hom}(\operatorname{NS}(T)/\mathrm{tors},\mathbf Z)\right).
\]
Thus \(N_1(T)\) is torsion-free and is separated by integral divisors.  The
completed monoid ring \(\Lambda_T\) is a domain.  Indeed, the effective
numerical monoid embeds in the torsion-free lattice \(N_1(T)\), hence is
cancellative and its ordinary monoid algebra is a domain.  An ample degree
is positive and proper on the effective numerical monoid, so every nonzero
completed series has a lowest nonzero homogeneous part, with only finitely
many monomials in that degree.  The lowest homogeneous part of a product is
the product of the two lowest homogeneous parts and is nonzero in the monoid
algebra.  Therefore \(\operatorname{Frac}(\Lambda_T)\) is defined.  Let
\(\mathcal K_T\) be an algebraic closure of this fraction field.  The loop
coordinate \(z\) remains independent: in particular, \(\mathcal K_T\)
contains no chosen root of \(z\).  We first isolate the cyclotomic part of
formal monodromy without choosing values for the exponentials of arbitrary
elements of \(\mathcal K_T\).

Let \(K\supset\C\) be an algebraically closed field of characteristic zero.
Choose a complement \(K=\C\oplus V\) of \(\mathbf Q\)-vector spaces, let
\(K[V]\) be the group algebra of the additive group of \(V\), and let
\(\Omega_V\) be an algebraic closure of its fraction field.  Thus
\(K\subset\Omega_V\).  The map
\[
 \operatorname{Exp}_V(c+v)=e^{2\pi i c}[v]
 \colon (K,+)\longrightarrow\Omega_V^\times
 \tag{5.0}
 \label{eq:universal-exponential}
\]
is a group homomorphism.  Since \(V\) is torsion-free, \(K[V]\) is a domain,
and its distinct monomials are linearly independent.  Consequently
\[
 \ker(\operatorname{Exp}_V)=\mathbf Z,
 \qquad
 \operatorname{Exp}_V(a)\text{ is a root of unity}
 \Longleftrightarrow a\in\mathbf Q.
 \tag{5.0a}
 \label{eq:universal-exponential-torsion}
\]

Apply Levelt--Turrittin after a ramification \(z=w^e\).  In the formal
decomposition, write the residue on a regular-singular block as
\(R=R_s+R_n\), with commuting semisimple and nilpotent parts.  In the
universal formal solution algebra, define a full turn of \(w\) on that block
by
\[
 M_{\mathrm{RS},V}
 =\operatorname{Exp}_V(R_s)\exp(2\pi iR_n),
\]
where \(\operatorname{Exp}_V\) is applied to the eigenvalues of \(R_s\).
The lift of one turn of the original \(z\)-disc acts on the formal symbols by
\(w\mapsto e^{2\pi i/e}w\), and on the blocks by their descent
intertwiners.  Its \(e\)-th power is \(M_{\mathrm{RS},V}\) on every returned
regular-singular block.  An integral change of residue representative
multiplies the formal solution by an integral power of \(w\), so it conjugates
this lifted action rather than changing it.  Since the coefficient and
descent matrices lie in \(K\subset\Omega_V\), this defines the framed
operator \(M_{\mathrm f,V}\) over \(\Omega_V\).  Ramified exponential factors
are thereby permuted rather than made single-valued by changing the frame.

\noindent\textbf{Choice independence.}
For every root of unity \(\zeta\in\C\), its algebraic multiplicity as an
eigenvalue of \(M_{\mathrm f,V}\) is independent of the complement \(V\), the
Levelt--Turrittin ramification and decomposition, and the algebraic closure
\(\Omega_V\).  It is also unchanged when the chosen algebraic closure of the
original coefficient field is replaced by a common algebraically closed
overfield.  When \(K=\C\), it is the usual multiplicity in the analytic
formal monodromy on the original \(z\)-disc.
We use here the functoriality and uniqueness of the ramified formal
decomposition.  That statement holds over any algebraically closed coefficient
field of characteristic zero, which is the algebraic form we need, and in that
generality it is the formal classification of differential modules over
\(K(\!(z)\!)\) \cite[Chapter~3]{PutSinger}.  Sabbah gives the
complex-analytic account, where the same decomposition appears alongside the
Riemann--Hilbert correspondence
\cite[Lecture~5, Section~5.c, proofs of Theorem~5.8 and
Proposition~5.10]{SabbahStokes}.
Consider one orbit of Levelt--Turrittin blocks under a turn of the original
disc, let \(d\) be its length, and write
\(A_i:E_i\longrightarrow E_{i+1}\) for its successive deck/descent maps,
with indices modulo \(d\).  Put
\(U=A_{d-1}\cdots A_0\) on \(E_0\).  A cyclic-block determinant gives
\[
 \chi_{\bigoplus_iE_i}(X)=\det(X^dI-U).
 \tag{5.0b}
 \label{eq:deck-orbit-polynomial}
\]
Put \(r=e/d\).  Equivariance of the descent datum gives the intrinsic return
relation
\[
 U^r=M_{\mathrm{RS},V}.
 \tag{5.0c}
 \label{eq:deck-return-power}
\]
Thus, if \(\mu\) is an eigenvalue of \(U\), then
\(\mu^r=\operatorname{Exp}_V(\rho)\) for a residue exponent \(\rho\), with
algebraic multiplicities read on the corresponding generalized eigenspaces.
If a root of unity \(\zeta\) occurs in
\eqref{eq:deck-orbit-polynomial}, then \(\mu=\zeta^d\) and
\(\operatorname{Exp}_V(\rho)=\zeta^e\); hence \(\rho\in\mathbf Q\) by
\eqref{eq:universal-exponential-torsion}.  Conversely, if \(\rho\) is
rational, then \(\mu^r\) is torsion, so every corresponding \(\mu\), and
hence every root of \(X^d-\mu\), is torsion.  On these rational residue
classes \(\operatorname{Exp}_V\) is the usual complex exponential; the
finite descent operator is the same intrinsic operator.  The exact roots of
unity and their algebraic multiplicities are therefore those of the ordinary
complex construction.  Nonrational residue classes contribute no root of
unity.  This criterion does not depend on \(V\) and does not require a
fractional residue factor to commute with the descent action.

For completeness, compare two Levelt--Turrittin choices after pulling both to
a common ramification \(z=\widetilde w^{\,E}\).  Equivariant uniqueness of the formal
decomposition identifies their total horizontal solution spaces and carries
the original-\(z\)-turn deck/descent operator for one choice to the scalar
extension of the other.  Thus the product of all orbit polynomials
\eqref{eq:deck-orbit-polynomial}, rather than any individually labelled
return map, is invariant under refinement and re-decomposition.  The same
argument after a common algebraically closed coefficient extension preserves
the intrinsic multiset of residue classes and finite descent characters.
The relation \eqref{eq:deck-return-power} is itself invariant under integral
changes of residue representative.  Hence the cyclotomic test above and its
algebraic multiplicities depend only on the original framed module, not on
the coefficient closure or complement.  Finally, for \(K=\C\) one has
\(V=0\) and
\(\operatorname{Exp}_V(a)=e^{2\pi i a}\), which gives the usual analytic
formal monodromy.

\begin{definition}
\label{def:framed-sixth-multiplicity}
Apply the preceding construction to the small even quantum connection over
\(\mathcal K_T((z))\), and put
\[
 \nu_6(T)=
 \operatorname{mult}_{e^{\pi i/3}}M_{\mathrm f,V}(T)+
 \operatorname{mult}_{e^{-\pi i/3}}M_{\mathrm f,V}(T).
\]
The multiplicities are algebraic multiplicities in the characteristic
polynomial.  The choice-independence argument above makes their sum independent
of all the auxiliary choices.
\end{definition}

The frame matters.  Adjoining roots of Novikov monomials is an algebraic
coefficient extension and is harmless by choice independence.  Adjoining a root of
\(z\), by contrast, would replace the original loop and could erase the deck
part of the operator.

\begin{lemmaOneA}[Frame transport over a coefficient field]
Let \(K\supset\C\) be an algebraically closed field of characteristic zero,
and let \(D,D'\) be differential modules over \(K(\!(z)\!)\).  Suppose
\(G\in\operatorname{GL}_n(K(\!(z)\!))\) is a gauge isomorphism from \(D\) to
\(D'\).  Then the framed operators \(M_{\mathrm f,V}(D)\) and
\(M_{\mathrm f,V}(D')\), defined using the original \(z\)-disc turn, are
conjugate after passage to a common coefficient extension.  In particular the
algebraic multiplicity of every root of unity is the same for \(D\) and
\(D'\).
\end{lemmaOneA}

\begin{proof}
After a common Levelt--Turrittin ramification, choose a formal solution matrix
\(Y\) for \(D\) in the formal solution algebra used in
Definition~\ref{def:framed-sixth-multiplicity}, and let \(\sigma\) denote one
turn of the original \(z\)-disc.  Since \(G\) has entries in
\(K(\!(z)\!)\), \(\sigma(G)=G\).  Hence
\[
 \sigma(GY)=G\sigma(Y).
\]
Thus \(GY\) is a solution frame for \(D'\) in which the original-disc turn has
the same matrix as in the frame \(Y\).  Comparing with any other solution frame
for \(D'\) conjugates that matrix by a constant matrix.  Therefore the framed
operators are conjugate and have the same characteristic polynomial.
\end{proof}

\begin{lemma}[Exact string and fixed-divisor shifts]
\label{lem:exact-low-degree-shifts}
Let a pullback of the even quantum connection of a smooth projective
variety $T$ be defined over a coefficient field containing its Novikov
coefficients.
\begin{enumerate}[label=(\roman*)]
\item Replacing a bulk point $\tau$ by $\tau+c\mathbf 1$ changes the
$z$-connection only by the scalar irregular twist $e^{-c/z}$.  Hence it does
not change the framed formal-monodromy eigenvalues on the original
$z$-disc, for any coefficient $c$ for which the shifted connection is
defined.
\item Let $D\in H^2(T)\otimes K$ be a bulk shift such that
$d\mapsto e^{\langle D,d\rangle}$ defines a multiplicative character on
the numerical Novikov image in the coefficient field.  Then the connection
at $\tau+D$ is obtained from that at $\tau$ by the coefficient
substitution
\[
 Q^d\longmapsto e^{\langle D,d\rangle}Q^d.
\]
This substitution fixes $\C$, so it preserves the algebraic
multiplicities of every root of unity in framed formal monodromy.
\end{enumerate}
\end{lemma}

\begin{proof}
Write the quantum connection in the convention of
\cite[(2.2)--(2.3) and Remark~2.3]{IritaniBlowup}.  The String Equation
implies that the quantum product is independent of the $H^0$ coordinate,
while the Euler field acquires $c\mathbf1$.  Thus the horizontal equation
$z^2\partial_zY=(U-z\mu)Y$ acquires the scalar term $cI$, and multiplication
of a solution by $e^{-c/z}$ gives the shifted equation.  This exponential is
single-valued under a full turn of the original $z$-disc and changes no
regular exponent or descent datum.

For a degree-two shift, the Divisor Equation says that the quantum product
depends on the divisor coordinate only through the combinations
$Q^d e^{\langle D,d\rangle}$; the degree-two coordinate does not occur in
the Euler field.  Therefore the full $z$-connection is the coefficient
base change displayed above.  The induced embedding of coefficient fields
fixes the complex roots of unity, and choice independence in
Definition~\ref{def:framed-sixth-multiplicity} gives the last assertion.
\end{proof}

\begin{remark}[Why a residual hypothesis remains]
\label{rem:residual-bulk-gap}
Flatness does give a normalized gauge for a genuinely formal bulk germ.
Writing the formal bulk variables collectively as $\eta$ and solving
\[
 \partial_{\eta_i}G=-z^{-1}(\phi_i\star_\eta)G,
 \qquad G|_{\eta=0}=I,
\]
recursively, the coefficient of total bulk degree $m$ has no power of $z$
below $-m$.  Thus every fixed bulk degree uses only integral powers of the
original loop coordinate and is harmless for framing.  After substituting a
reconstruction coordinate containing Novikov coefficients, however, the sum
over all bulk degrees need not have a uniform lower Laurent bound in $z$.
We therefore do not promote this degreewise gauge to an element of a Laurent
series field and do not use it to identify intrinsic framed monodromy at a
reconstruction value.  Hypotheses~5.7R and~5.7T below isolate exactly the two
places where such an identification is still required.  The rank-two theorem
below avoids this gauge altogether on the formal germ where it applies.
\end{remark}

\begin{lemma}[Numerical Novikov base change]
\label{lem:numerical-base-change}
The quotient from the effective homology monoid to its image in \(N_1(T)\)
extends continuously to completed monoid rings, and the quantum connection
and the two comparison theorems base-change along it.
\end{lemma}

\begin{proof}
Filter both monoids by intersection with an ample divisor.  Effective homology
classes below a fixed cutoff form a finite set, as follows from the finite-type
Chow varieties of cycles of bounded degree.  The quotient is therefore defined
on every truncated monoid ring and passes to the inverse limit.  Explicitly,
the coefficient of \(Q^{\bar d}\) in the image of \(\sum_da_dQ^d\) is
\[
 \sum_{d\mapsto\bar d}a_d.
\]
The sum is finite at every cutoff, and additivity of the quotient makes this
coefficientwise pushforward a unital homomorphism for completed convolution.
Thus the quantum products descend by grouping Gromov--Witten coefficients with
the same numerical class.  Pushforward of curves, pairing with
\(c_1(N_{C/T})\), and the derivations \(Q\partial_Q\) are constant on these
fibres.  At each cutoff, apply the resulting ring homomorphism to the
comparison map and to its stated inverse.  Scalar extension preserves their
two inverse identities.  Grouping coefficients identifies the scalar-extended
quantum module with the numerical quantum module, and the preceding
fiberwise constancy identifies its connection operators and divisor
substitutions.  Passing through the compatible cutoffs proves the asserted
completed base change.
\end{proof}

\begin{proposition}[Formal-germ rigidity of an isolated rank-two packet]
\label{prop:ranktwo-framed-germ}
Let $K$ be an algebraically closed coefficient field of characteristic zero
and let the even quantum connection be pulled back to a formal smooth bulk
germ $B=K[[t_1,\ldots,t_m]]$.  Suppose that at the closed point the Euler
operator has an eigenvalue whose generalized eigenspace has rank two, its
nilpotent part is a single nonzero Jordan block, and every other Euler
eigenvalue differs from it by a unit of $K$.  Then the corresponding
rank-two spectral cluster remains a single nonzero Jordan block on the
formal germ.  After scalar centering and the canonical elementary
modification of Section~\ref{sec:atomic-one-step}, its residue $R$ has
constant characteristic polynomial on $B$.  Consequently the algebraic
multiplicities of the primitive sixth roots contributed by this block are
constant on the formal germ.
\end{proposition}

\begin{proof}
The unit separation from the other Euler eigenvalues gives a formal spectral
projector onto the rank-two cluster.  Put
$U_0=\lambda I+N$ on this cluster, with $\operatorname{tr}N=0$, and let
$C_a$ denote quantum multiplication by the bulk tangent vector
$\partial_{t_a}$.  Flatness gives the coefficient identities
\[
 [U,C_a]=0,\qquad
 \partial_{t_a}U=C_a+[C_a,\mu].
\]
Let $P$ be the formal spectral projector onto this cluster and use the
induced connection $D_a$ on $P(H)$.  For an endomorphism preserving the
cluster, $D_a$ is the projected derivative, so
\[
 D_aU_0=P(\partial_{t_a}U)P.
\]
Because $[U,C_a]=0$, the operator $C_a$ commutes with $P$.  Compressing the
second flatness identity therefore gives
\[
 D_aU_0=C_{a,0}+[C_{a,0},\mu_0],
 \qquad C_{a,0}=PC_aP,\quad \mu_0=P\mu P.
\]
Because the reduction of $U_0$ at the closed point is nonscalar, $U_0$ is a
regular $2\times2$ matrix over the local ring $B$.  After choosing the
closed-point Jordan frame, one off-diagonal entry of $N$ is a unit; solving
$[N,X]=0$ directly in that frame gives the commutant $B I\oplus B N$.
Hence
\[
 C_{a,0}=p_aI+q_aN.
\]
Taking traces in the block evolution gives $p_a=\partial_{t_a}\lambda$ and
therefore
\[
 D_aN=q_aN+q_a[N,\mu_0].
\]
Set $d=\det N$.  Since determinant is invariant under change of frame,
its ordinary derivative can be computed with the covariant derivative
$D_a$.  For a traceless $2\times2$ matrix,
$\operatorname{adj}(N)=-N$ and $N^2=-dI$.  Cyclicity of trace then gives
\[
 \partial_{t_a}d
 =\operatorname{tr}(\operatorname{adj}(N)D_aN)
 =2q_a d.
\]
Since $d$ vanishes at the closed point, a lowest-homogeneous-term argument
in $K[[t_1,\ldots,t_m]]$ forces $d=0$.  Thus $N^2=0$ everywhere.  A matrix
entry of $N$ is a unit at the closed point, so $N$ remains nonzero and its
image is a line subbundle.

The proof of Proposition~\ref{prop:rank2-rigidity} is purely algebraic: it
uses only flatness, Frobenius self-adjointness, anti-self-adjointness of the
grading operator, and the nondegenerate Poincar\'e pairing on the separated
spectral factor.  Those identities hold coefficientwise over the present
formal germ, so the same elementary modification is regular in every base
direction and its residue satisfies
\[
 D_aR=[G_a,R].
\]
Hence $\operatorname{tr}R$ and $\det R$, and therefore the characteristic
polynomial of $R$, are constant.  By
Proposition~\ref{prop:residue-discriminant-exponents}, the residue
eigenvalues are representatives of the formal exponent classes; the
integral shear changes them only by integers.  The eigenvalues of the
regular formal monodromy are therefore constant, including in resonant
cases where only the unipotent part could vary.  This proves the cyclotomic
multiplicity statement.
\end{proof}

\begin{remark}[Formal-germ scope]
\label{rem:ranktwo-germ-scope}
Proposition~\ref{prop:ranktwo-framed-germ} is deliberately a theorem over a
formal bulk germ with a fixed coefficient field.  It does not identify an
Iritani or Iritani--Koto reconstruction value with a point of that germ.
The comparison theorems are formulated over graded Laurent completions in
Novikov variables, whereas the proposition deforms genuinely formal bulk
variables over $K$.  Substituting a Laurent-Novikov reconstruction expression
for those formal variables is an additional operation, and no theorem
legitimizing that substitution for framed monodromy is assumed here.  Thus
the proposition proves the local rank-two mechanism behind
Hypothesis~5.7R; it is not a proof of that hypothesis at a reconstruction
value.
\end{remark}

\begin{hypothesisR}[Reconstruction-tail invariance]
Consider only the residual reconstruction displacements that occur in the
blowup and projective-bundle decompositions used in
Proposition~\ref{prop:framed-operations}: the $O(u)+s_j$ tails on the target
summands and the corresponding residual ambient shift, after the scalar
$H^0$ term and every fixed $H^2$ term whose Novikov character descends to
the relevant coefficient ring have been removed by
Lemma~\ref{lem:exact-low-degree-shifts}.  After the coefficient
specialization displayed in the construction, replacing zero bulk by one of
these residual reconstruction displacements does not change the algebraic
multiplicities of $e^{\pi i/3}$ and $e^{-\pi i/3}$ in framed formal
monodromy on the original $z$-disc.

No invariance under an arbitrary bulk displacement is asserted.
\end{hypothesisR}



\begin{proposition}[Framed operation formulas, conditional]
\label{prop:framed-operations}
Assume Hypothesis~5.7R.
Let \(V\to T\) be a vector bundle of rank \(r\ge2\).  Then
\[
 \nu_6(\PP_T(V))=r\nu_6(T).
 \tag{5.2}
 \label{eq:projective-bundle-nu}
\]
Let \(\widetilde T=\operatorname{Bl}_C T\), where \(C\subset T\) is smooth of
codimension \(c\ge2\).  In Iritani's comparison field one has
\[
 \nu_6(\widetilde T)=\nu_6(T)+
 \sum_{j=0}^{c-2}\nu_6(C;\chi_j),
 \tag{5.3}
 \label{eq:blowup-nu}
\]
where \(\chi_j\) is the numerical Novikov specialization of the \(j\)-th
center summand.  The endpoint terms in \eqref{eq:blowup-nu} are the intrinsic small
invariants; the center terms are computed after the displayed specializations.
\end{proposition}

\begin{proof}
We first treat the blowup.  Iritani's Theorem~5.18 gives a formal coordinate
isomorphism and a quantum-$D$-module decomposition
\[
 \operatorname{QDM}(\widetilde T)\cong
 \operatorname{QDM}(T)\oplus
 \operatorname{QDM}(C)^{\oplus(c-1)}
\]
over a Laurent extension in the exceptional Novikov variable
\cite[Theorem~5.18]{IritaniBlowup}.  The comparison map and its inverse use
only integral powers of the original loop coordinate $z$.  After passing the
coefficient ring to an algebraically closed Laurent coefficient field,
Lemma~5.1A therefore transports framed formal monodromy across the comparison
without ramifying the $z$-disc.

For a center summand put $u=q^{-1/(c-1)}$.  Iritani's reconstruction
coordinates have the initial form
\[
 \varsigma_j=\varsigma_j^\circ+s_j,
 \qquad
 \varsigma_j^\circ=-(c-1)\lambda_j+h_{C,j}+O(u),
 \tag{5.4}
 \label{eq:blowup-reconstruction-coordinate}
\]
with $\lambda_j=\zeta_j u^{-1}$
\cite[Section~5.8.2 and (5.27)--(5.30)]{IritaniBlowup}.  The associated
numerical Novikov specialization is
\[
 \chi_j(Q_C^d)=Q^{i_*d}u^{\rho_C\cdot d}.
 \tag{5.4a}
 \label{eq:center-novikov-specialization}
\]
The unnormalized connection is naturally read after the Laurent localization
containing $u^{-1}$; in particular we make no claim that the term
$\lambda_j$ belongs to a nonlocalized center ring.  Lemma~\ref{lem:exact-low-degree-shifts}
removes the scalar $H^0$ term $-(c-1)\lambda_j$ and the fixed divisor
$h_{C,j}$ exactly.  The latter character descends through
\eqref{eq:center-novikov-specialization} because $h_{C,j}$ is a scalar
multiple of $\rho_C$.  The remaining $O(u)+s_j$ reconstruction tail is the
center instance covered by Hypothesis~5.7R.  The normalization of the
logarithmic coordinate is fixed in \cite[Section~5.8.1]{IritaniBlowup}; a
residual constant logarithmic character would again be a fixed divisor
shift and hence is already covered by Lemma~\ref{lem:exact-low-degree-shifts}.

The ambient summand is analogous.  Its reconstruction coordinate is
$\tau=\tau^\circ+t$ with $\tau^\circ=O(q^{-1})$ in the Laurent expansion of
Theorem~5.18.  Hypothesis~5.7R is invoked only for this residual ambient
reconstruction tail.  Thus, after the exact low-degree normalizations, the
ambient contribution has the intrinsic multiplicity of $T$, while the
$j$-th center contribution has the intrinsic multiplicity of the small
connection of $C$ after the specialization $\chi_j$.  Direct-sum additivity
gives \eqref{eq:blowup-nu}.

The coefficient extensions do not alter the endpoint invariants.  On
numerical curve lattices the blowup map
\[
 \bar\beta\longmapsto
 \bigl(\phi_*\bar\beta,-E\cdot\bar\beta\bigr)
\]
is injective because
$N^1(\widetilde T)=\phi^*N^1(T)\oplus\mathbf Z[E]$.  Hence the numerical
Novikov field of $\widetilde T$ embeds in the comparison field.  The center
map \eqref{eq:center-novikov-specialization} need not be injective, which is
why its contribution is denoted $\nu_6(C;\chi_j)$ rather than $\nu_6(C)$.

For a projective bundle, Iritani--Koto's Theorem~5.1 gives the corresponding
decomposition into $r$ copies of the quantum $D$-module of $T$ over a
Laurent coefficient extension \cite[Theorem~5.1]{IritaniKoto}.  Again the
comparison and its inverse use only integral powers of the same loop
coordinate, so Lemma~5.1A applies.  With $u=q^{-1/r}$ their reconstruction
coordinates satisfy
\[
 \varsigma_j=\varsigma_j^\circ+s_j,
 \qquad
 \varsigma_j^\circ
 =r\lambda_j-\frac{2\pi\mathrm{i}j}{r}c_1(V)+O(u)
\]
\cite[Theorem~5.1(4), Proposition~5.6]{IritaniKoto}.  Lemma~\ref{lem:exact-low-degree-shifts}
removes the scalar term and the fixed divisor term.  The initial-coordinate
normalization is fixed by \cite[(5.11)--(5.12)]{IritaniKoto}.  Hypothesis~5.7R
is then used only for the residual $O(u)+s_j$ reconstruction tail.  The base
Novikov map is injective \cite[(1.1) and Remark~5.2]{IritaniKoto}; hence every
summand contributes the intrinsic value $\nu_6(T)$, and direct-sum additivity
gives \eqref{eq:projective-bundle-nu}.

The projective-bundle theorem is stated after a global-generation
normalization.  Tensoring $V$ by a sufficiently negative line bundle does
not change $\PP_T(V)$, and the intrinsic Novikov embedding is unchanged by
this twist \cite[Remark~1.2 and Remark~5.2]{IritaniKoto}.  Finally, both
comparison decompositions preserve cohomological parity on the even
restriction used here.  Iritani records this explicitly for blowups
\cite[Section~2]{IritaniNotes}; for projective bundles it follows directly
from the pullback, pushforward, Fourier-transform, and characteristic-class
operations in Theorem~5.1 after odd bulk variables are set to zero.  The
adjoined Novikov roots are coefficient variables, not roots of $z$.
\end{proof}

The center map in \eqref{eq:blowup-nu} is not necessarily injective
\cite[(5.15)]{IritaniBlowup}.  Thus generic vanishing for \(C\) does not by
itself imply that its specialized term vanishes.  The next lemma supplies the
missing implication.

\begin{definition}
\label{def:strict-novikov-admissible}
A monomial map \(\chi:\Lambda_T\to A\) is \emph{strictly
Novikov-admissible} if \(A\) is a complete separated valued \(\C\)-domain
with domain associated graded, every monomial has nonzero image, and
\[
 v(\chi(Q^d))=\ell(d),
 \tag{5.5}
 \label{eq:strict-novikov-valuation}
\]
where \(\ell\) is positive and proper on the effective numerical monoid.
\end{definition}

We write \(\nu_6(T;\chi)\) for the framed small-even invariant after scalar
extension along such a map.

The center maps in \eqref{eq:blowup-nu} are strictly admissible.  Give the exceptional
Laurent variable valuation zero and use an ample divisor \(H\) on \(T\).  Then
\[
 v_H\!\left(Q^{i_*d}q^{-\rho_C\cdot d/(c-1)}\right)
 =(H|_C)\cdot d.
\]
The restriction \(H|_C\) is ample, which gives positivity and properness; the
associated graded Laurent monoid ring is a domain.

\begin{proposition}[Direct specialized low-dimensional vanishing]
\label{prop:direct-specialized-lowdim}
Let \(\chi:\Lambda_T\to A\) be strictly Novikov-admissible.
\begin{enumerate}[label=(\roman*)]
\item If the canonical divisor \(K_T\) is nef, then
\[
  \nu_6(T;\chi)=0
\]
unconditionally.
\item If \(T=\PP^1\) or \(T=\PP^2\), then
\[
  \nu_6(T;\chi)=0
\]
unconditionally.
\item Let \(T=\PP_C(V)\) be a geometrically ruled surface with
\(g(C)\ge1\).  Assuming Hypothesis~5.7R,
\[
  \nu_6(T;\chi)=0.
\]
\end{enumerate}
\end{proposition}

\begin{proof}
For (i), work after scalar extension to an algebraic closure of
\(\operatorname{Frac}A\).  At small bulk write
\[
 \nabla_{z\partial_z}=z\partial_z-z^{-1}U+\operatorname{Gr},
 \qquad
 U=\sum_{\beta}\chi(Q^\beta)U_\beta,
\]
where \(U_\beta\) is the coefficient of Euler multiplication in numerical
class \(\beta\).  Put
\[
 g=\operatorname{Gr}+\frac12E,
\]
with \(E\) equal to the identity on cohomological degrees of parity
opposite to \(\dim T\) and zero on the other parity.  The operator \(g\) has
integral eigenvalues.  The genus-zero dimension constraint gives
\[
 U_\beta:H^i(T)\longrightarrow
 H^{i+2-2c_1(T)\cdot\beta}(T),
\]
so, because quantum multiplication preserves parity, \(U_\beta\) raises
\(g\)-weight by
\[
 m_\beta=1-c_1(T)\cdot\beta.
\]
If \(K_T\) is nef, then \(m_\beta\ge1\) for every effective
\(\beta\).  Make the half-parity correction explicitly by replacing the
\(z\)-connection with
\[
 \nabla^{\mathrm{par}}_{z\partial_z}
 =\nabla_{z\partial_z}+\frac12E
 =z\partial_z-z^{-1}U+g.
\]
With the gauge convention \(G\nabla G^{-1}\), take \(G=z^g\).  Since
\([g,U_\beta]=m_\beta U_\beta\), one has
\[
 z^gU_\beta z^{-g}=z^{m_\beta}U_\beta,
 \qquad
 z^g(z\partial_z)z^{-g}=z\partial_z-g,
\]
and therefore
\[
 z^g\nabla^{\mathrm{par}}_{z\partial_z}z^{-g}
 =z\partial_z-
   \sum_\beta \chi(Q^\beta)z^{m_\beta-1}U_\beta.
\]
Thus the specialized connection is regular singular.  This conclusion is
unchanged when the noninjective map \(\chi\) groups several numerical
classes: every grouped summand still has nonnegative \(z\)-order.
The residue consists only of the terms with \(m_\beta=1\), equivalently
\(c_1(T)\cdot\beta=0\).  Each such operator raises the finite
\(g\)-filtration by one, so their sum is nilpotent.  The parity-corrected
regular monodromy is therefore unipotent.  Undoing the half-parity correction
permits only eigenvalues \(1\) and \(-1\), because the transformed residue
preserves parity.  Hence no primitive sixth root occurs.  Strict
Novikov-admissibility is used here only to ensure that the specialized small
connection is a well-defined coefficientwise specialization; injectivity is
not used.

For (ii), let \(\ell\) be the line class and put
\(a=\chi(Q^\ell)\).  Strict admissibility gives \(a\ne0\).  Over an
algebraic closure of \(\operatorname{Frac}A\), the specialized small
quantum relation for \(\PP^m\), \(m=1,2\), is
\[
 H^{m+1}=a.
\]
Euler multiplication by \((m+1)H\) therefore has the \(m+1\) distinct
eigenvalues
\[
 (m+1)\zeta a^{1/(m+1)},
 \qquad \zeta^{m+1}=1.
\]
Every spectral block has rank one, so
Lemma~\ref{lem:simple-euler-block} gives trivial framed regular monodromy
on every block.

For (iii), let \(f\) be the fibre class.  Every map from a connected
nodal genus-zero curve to \(C\) is constant: it is constant on every
irreducible rational component, and connectedness forces the same image
point on all components.  Hence every nonzero curve class contributing to
the genus-zero quantum product of \(T\) is a positive multiple of \(f\).
The small quantum differential module is therefore obtained by scalar
extension from a one-variable module over \(\C[[Q^f]]\).  The restriction
of \(\chi\) to the fibre ray is injective, because
\[
 v(\chi(Q^{nf}))=n\ell(f),
 \qquad \ell(f)>0,
\]
so distinct powers have distinct valuation.  On the other hand,
Hypothesis~5.7R and Proposition~\ref{prop:framed-operations} give the
intrinsic projective-bundle identity
\[
 \nu_6(T)=2\nu_6(C)=0,
\]
where the last equality follows by applying (i) to the identity Novikov
specialization of \(C\), since \(K_C\) is nef for \(g(C)\ge1\).
The intrinsic small module and the \(\chi\)-specialized module are both
scalar extensions of this same one-variable differential module.  More
explicitly, let
\[
 K_f=\operatorname{Frac}(\C[[Q^f]]),
 \qquad
 K_{f,\chi}=\operatorname{Frac}(\chi(\C[[Q^f]]))
 \subset\operatorname{Frac}(A).
\]
The injective restriction of \(\chi\) identifies \(K_f\) with
\(K_{f,\chi}\).  Via this identification, choose a common algebraically
closed overfield \(\Omega\) containing \(K_f\) and
\(\operatorname{Frac}(A)\).  After scalar extension to \(\Omega\), the
intrinsic one-variable module and its \(\chi\)-specialization become
literally the same differential module.  Definition~\ref{def:framed-sixth-multiplicity} and
its choice-independence under common coefficient-field extension therefore
give identical root-of-unity multiplicities.  The specialized value also
vanishes.
\end{proof}

\begin{remark}[Where the tagging hypothesis is still used]
\label{rem:tagging-scope}
Proposition~\ref{prop:direct-specialized-lowdim} shows that
Hypothesis~5.7T is not needed for nef-canonical centers, for \(\PP^1\) or
\(\PP^2\), or---assuming only Hypothesis~5.7R---for ruled surfaces over
positive-genus curves.  In the present proof of four-dimensional birational
invariance, its substantive use is confined to minimal rational ruled
surfaces and to nonminimal surfaces.  Eliminating it there would require a
separate support/base-change statement for applying the blowup or
projective-bundle comparison after the external center specialization; no
such statement is assumed here.
\end{remark}

\begin{hypothesisT}[Divisor-tagging specialization invariance]
Let $\chi:\Lambda_T\to A$ be strictly Novikov-admissible.  Choose integral
divisors $D_1,\ldots,D_\rho$ whose pairings separate $N_1(T)$ and define
\[
 \chi_{\boldsymbol t}(Q^d)
 =\chi(Q^d)\exp\left(\sum_i t_i(D_i\cdot d)\right).
\]
After passing to algebraic closures of the relevant fraction fields, the
intrinsic framed operator of the $\chi$-specialized small connection and
the tagged module have the same algebraic multiplicities at
$e^{\pi i/3}$ and $e^{-\pi i/3}$.

This hypothesis concerns only the specialization step in this tagging
family.  It does not assert reconstruction invariance or invariance under
an arbitrary divisor-valued bulk deformation.
\end{hypothesisT}

\begin{lemma}[Divisor tagging, conditional]
\label{lem:divisor-tagging}
Assume Hypothesis~5.7T.
Let \(\chi:\Lambda_T\to A\) be strictly Novikov-admissible.  If the intrinsic
small even quantum connection of \(T\) has no primitive-sixth framed
monodromy, neither does its specialization along \(\chi\).
\end{lemma}

\begin{proof}
If \(N_1(T)=0\), then the numerical effective monoid is \(\{0\}\), so the
coefficient map is already injective and there is nothing to tag.  Assume
otherwise.
Choose integral divisors \(D_1,\ldots,D_\rho\) whose pairings separate
\(N_1(T)\), and introduce formal variables
\(\boldsymbol t=(t_1,\ldots,t_\rho)\).  Define
\[
 \chi_{\boldsymbol t}(Q^d)=\chi(Q^d)
 \exp\!\left(\sum_i t_i(D_i\cdot d)\right).
 \tag{5.6}
 \label{eq:divisor-tagging-map}
\]
This formula extends continuously from monomials to the completed Novikov
ring.  Indeed, strict Novikov admissibility gives
\(v(\chi(Q^d))=\ell(d)\), with \(\ell\) positive and proper on the
effective numerical monoid.  Hence only finitely many classes contribute
below any valuation cutoff, so the tagged sum converges coefficientwise in
the complete ring \(A[[\boldsymbol t]]\) and defines a continuous unital
homomorphism.  We now prove that it is injective.  Take
\(f=\sum_d c_dQ^d\ne0\) and put
\[
 \mu=\min\{v(\chi(Q^d)):c_d\ne0\}
     =\min\{\ell(d):c_d\ne0\},
 \qquad
 S_\mu=\{d:c_d\ne0,\ \ell(d)=\mu\}.
 \label{eq:divisor-tagging-initial-support}
\]
Properness in \eqref{eq:strict-novikov-valuation} makes \(S_\mu\) finite
and nonempty.  Write \(\operatorname{in}_v\) for the initial class in the
associated graded.  The degree-\(\mu\) initial form is
\[
 \operatorname{in}_\mu\!\left(\chi_{\boldsymbol t}(f)\right)
 =\sum_{d\in S_\mu}c_d\operatorname{in}_v\!\left(\chi(Q^d)\right)
 \exp\!\left(\sum_i t_i(D_i\cdot d)\right)
 \ \in\ \operatorname{gr}_v(A)[[\boldsymbol t]],
 \label{eq:divisor-tagging-initial-form}
\]
and every displayed summand coefficient is nonzero.  Every term outside \(S_\mu\) has valuation
strictly greater than \(\mu\), so a nonzero initial form proves
\(\chi_{\boldsymbol t}(f)\ne0\).  This is a finite combination of distinct
exponential characters.  The finitely many exponent vectors are distinct
integral vectors.  Choose \(a=(a_i)\in\Z^\rho\) outside the finitely many
integral hyperplanes on which two dot products agree, and substitute
\(t_i=a_i s\).  The resulting exponents are distinct integers.  Their
derivatives at \(s=0\) form a Vandermonde matrix with nonzero integer
determinant, which stays invertible over
\(\operatorname{Frac}(\operatorname{gr}_v A)\) because that field has
characteristic zero.  The characters are therefore linearly independent, and
the initial form displayed above is nonzero.

The domain hypothesis on \(\operatorname{gr}_v(A)\) is used twice, and only
twice.  It provides the field \(\operatorname{Frac}(\operatorname{gr}_v A)\)
of the Vandermonde step just made, and it makes \(v\) multiplicative on
\(A\), since initial forms of nonzero elements then have nonzero product.  The
valuation of a monomial image is \(\ell(d)\) by hypothesis
\eqref{eq:strict-novikov-valuation}; multiplicativity is what preserves that
value after multiplication by the tag, a unit of valuation zero.  Terms with
\(\ell(d)>\mu\) therefore contribute nothing in degree \(\mu\), and each
\(d\in S_\mu\) contributes its initial form with \(c_d\) intact.  Were the
associated graded to have zero divisors, \(v\) would be a filtration rather
than a valuation, and the displayed degree-\(\mu\) identity could fail.

Injectivity of
\eqref{eq:divisor-tagging-map} embeds the intrinsic numerical Novikov
fraction field into the fraction field of $A[[\boldsymbol t]]$, fixing
$\C$.  Therefore the module obtained by the tagged map is a scalar extension
of the intrinsic small module, and their framed characteristic polynomials
have the same primitive-sixth multiplicities after passage to a common
algebraic closure.

On the other hand, setting $\boldsymbol t=0$ in the tagged family gives the
$\chi$-specialized small connection.  Hypothesis~5.7T is precisely the
remaining assertion that this specialization does not change the two
primitive-sixth multiplicities.  Combining these two statements proves the
lemma.  Notice that no pro-Laurent or Hahn-field framed operator is invoked:
the hypothesis is attached directly to the specialization step for which it
is needed.
\end{proof}

\begin{proposition}[Low-dimensional vanishing, conditional]
\label{prop:low-dimensional-vanishing}
Assume Hypotheses~5.7R and~5.7T.
Let \(T\) be a point, a smooth projective curve, or a smooth projective
surface.  For every strictly Novikov-admissible specialization \(\chi\),
\[
 \nu_6(T;\chi)=0.
 \tag{5.7}
 \label{eq:low-dimensional-vanishing}
\]
\end{proposition}

\begin{proof}
The direct specialized cases are already covered by
Proposition~\ref{prop:direct-specialized-lowdim}: this includes points and
curves of positive genus through the nef-canonical case, the two projective
spaces \(\PP^1,\PP^2\), every minimal surface with nef canonical bundle,
and, assuming Hypothesis~5.7R, ruled surfaces over positive-genus curves.

It remains to treat minimal rational ruled surfaces and nonminimal surfaces.
First work intrinsically, before the external specialization.  If
\(T=\PP_{\PP^1}(V)\) is a minimal rational ruled surface, then
Proposition~\ref{prop:framed-operations} and the unconditional value
\(\nu_6(\PP^1)=0\) give \(\nu_6(T)=0\).  Every nonminimal smooth
projective surface is obtained from a minimal surface by point blowups; under
Hypothesis~5.7R the blowup formula
\eqref{eq:blowup-nu} adds only point terms, and those are zero.  Thus the
intrinsic small invariant vanishes for the remaining surfaces as well.

For these remaining rational ruled and nonminimal cases we now use
Hypothesis~5.7T through Lemma~\ref{lem:divisor-tagging}.  It carries the
intrinsic vanishing to the strictly Novikov-admissible specialization
\(\chi\), proving \eqref{eq:low-dimensional-vanishing}.  No use of
Hypothesis~5.7T is made in the direct cases of
Proposition~\ref{prop:direct-specialized-lowdim}.
\end{proof}

\begin{theorem}[Conditional birational invariance through dimension four]
\label{thm:nu6-birational-invariance}
Assume Hypotheses~5.7R and~5.7T.
If \(Y\) and \(Y'\) are birational smooth projective complex varieties of
the same dimension at most four, then
\[
 \nu_6(Y)=\nu_6(Y').
 \tag{5.8}
 \label{eq:nu-birational-invariance}
\]
Consequently, for smooth projective threefolds,
\[
 Y\times\PP^1\dashrightarrow Y'\times\PP^1
 \quad\Longrightarrow\quad \nu_6(Y)=\nu_6(Y').
\]
\end{theorem}

\begin{proof}
Weak factorization decomposes a birational map between smooth varieties into
blowups and blowdowns with smooth centers of codimension at least two
\cite{AKMW}.  In dimension at most four every nontrivial center has dimension
at most two.  The calculation following
Definition~\ref{def:strict-novikov-admissible} shows that every center
specialization in \eqref{eq:blowup-nu} is strictly admissible.
Proposition~\ref{prop:low-dimensional-vanishing} therefore makes every center
term zero, so the endpoint values agree at each step.

If \(Y,Y'\) are smooth threefolds and
\(Y\times\PP^1\dashrightarrow Y'\times\PP^1\), apply the first assertion in
dimension four and \eqref{eq:projective-bundle-nu}:
\[
 2\nu_6(Y)=\nu_6(Y\times\PP^1)
 =\nu_6(Y'\times\PP^1)=2\nu_6(Y').
\]
Cancelling two in \(\Z\) proves the stabilization assertion.
\end{proof}

\begin{lemma}[Multiplicity-one Euler blocks]\label{lem:simple-euler-block}
Let $T$ be smooth projective and consider its even quantum connection at an
even bulk point over an algebraically closed coefficient field.  If a
generalized eigenspace of Euler multiplication has rank one, that block
contributes only the eigenvalue $1$ to framed regular monodromy.  In
particular, at the small point it contributes nothing to $\nu_6(T)$.
\end{lemma}

\begin{proof}
Write the horizontal system at the chosen even bulk point in the standard form
\[
 z^2\partial_zY=(U-z\mu)Y,
\]
where $U$ is Euler quantum multiplication and $\mu$ is the grading operator.
The Frobenius property makes $U$ self-adjoint for the Poincar\'e pairing,
while $\mu$ is anti-self-adjoint.  Generalized eigenspaces belonging to
distinct eigenvalues of the self-adjoint operator $U$ are orthogonal, so the
pairing restricts nondegenerately to the one-dimensional block.  If $P$ is
its spectral projector, anti-self-adjointness on this nondegenerate line
therefore gives
\[
 P\mu P=0.
\]

Formally split this line from the other Euler eigenspaces by the normalized
gauge $g=I+zg_1+z^2g_2+\cdots$ whose positive coefficients are block-off-
diagonal.  The coefficient of $z$ on the resulting scalar block is
\[
 -P\mu P+P[U,g_1]P=0;
\]
the logarithmic-derivative term of the gauge begins in order $z^2$.
Consequently the scalar equation has the form
\[
 z^2\partial_zy=(u+z^2h(z))y,
 \qquad h(z)\in K[[z]],
\]
for the corresponding Euler eigenvalue $u$.  Its normalized formal solution
is $e^{-u/z}v(z)$ with $v(z)\in1+zK[[z]]$.  No fractional power of the
original loop coordinate and no logarithm occurs, so the framed regular
monodromy on this block is the identity.
\end{proof}

\begin{proposition}[Products with projective space]
\label{prop:projective-product-nu}
For every smooth projective variety $T$ and every $m\ge0$,
\[
 \nu_6(T\times\PP^m)=(m+1)\nu_6(T).
\]
This statement is unconditional.
\end{proposition}

\begin{proof}
The product formula for Gromov--Witten theory
\cite{BehrendProduct}, followed by the numerical Novikov base change of
Lemma~\ref{lem:numerical-base-change}, identifies the small numerical
quantum connection of $T\times\PP^m$ with the tensor-product quantum
connection of the two factors.  On even cohomology this is simply
$H^{\rm ev}(T)\otimes H^{\rm ev}(\PP^m)$, since projective space has no odd
cohomology.  After adjoining $q^{1/(m+1)}$ in the projective-space Novikov
coefficient field, quantum multiplication by
$c_1(\PP^m)=(m+1)H$ has the $m+1$ distinct eigenvalues
\[
 (m+1)\zeta q^{1/(m+1)},\qquad \zeta^{m+1}=1.
\]
By Lemma~\ref{lem:simple-euler-block}, each corresponding rank-one formal
factor has trivial framed regular monodromy.  The Levelt--Turrittin
formal decomposition is compatible with tensor products of differential
modules.  Hence tensoring a formal factor of the quantum differential module
of $T$ with one of these rank-one projective-space factors only adds its
scalar irregular exponential factor and leaves the original-$z$-disc
regular framed operator of the $T$ factor unchanged.  The extension adjoining
$q^{1/(m+1)}$ is in the coefficient field and is fixed by the $z$-turn.
Consequently the product module has $m+1$ copies of every framed cyclotomic
factor of $T$, counted with algebraic multiplicity.  This gives the formula.
\end{proof}

\begin{corollary}\label{cor:p3-nu6}
The small even framed invariant of projective three-space satisfies
\[
 \nu_6(\PP^3)=0.
\]
\end{corollary}

\begin{proof}
Apply Proposition~\ref{prop:projective-product-nu} with $T=\operatorname{pt}$
and $m=3$.
\end{proof}

\begin{proposition}[The cubic packet]
\label{prop:cubic-packet}
For every smooth cubic threefold \(X\),
\[
 \nu_6(X)=2.
 \tag{5.9}
 \label{eq:cubic-packet-value}
\]
\end{proposition}

This proposition is unconditional and uses neither Hypothesis~5.7R nor
Hypothesis~5.7T.

\begin{proof}
Use the notation of Proposition~\ref{prop:cubic-block-data} and put
\(F_X=\overline{\operatorname{Frac}\C[[q]]}\).  That
proposition separates the small even connection into two scalar blocks and
the rank-two zero block \eqref{eq:cubic-integral-block-reduction}, using
only an algebraic Novikov extension and integral powers of the original
loop coordinate.  We treat the rank-two block first and return to the scalar
blocks below.

For the rank-two block, seek a formal solution in the form
\[
 \widetilde S_3=z^\rho(1+O(z)),
 \qquad
 \widetilde S_4=cz^{\rho+1}(1+O(z)).
 \tag{5.9f}
 \label{eq:cubic-indicial-ansatz}
\]
The coefficient of \(z^{\rho+1}\) in the first coordinate of
\eqref{eq:cubic-integral-block-reduction}, and that of \(z^{\rho+2}\) in
the second, give respectively
\[
 \rho=2c-\frac{19}{18},
 \qquad
 c(\rho+1)=\frac{19}{18}c-\frac8{81}.
 \tag{5.9g}
 \label{eq:cubic-indicial-coefficients}
\]
Eliminating \(c\) gives
\[
 \frac12\left(\rho+\frac{19}{18}\right)
 \left(\rho-\frac1{18}\right)=-\frac8{81},
 \qquad\text{hence}\qquad
 \rho^2+\rho+\frac5{36}=0.
 \tag{5.9h}
 \label{eq:cubic-indicial-polynomial}
\]
This is the coefficient comparison in Cai's displayed rank-two
ansatz~\cite[p.~6]{Cai}.

The quadratic formula applied to \eqref{eq:cubic-indicial-polynomial} gives
\[
 \rho=-\frac16,\ -\frac56.
 \tag{5.9i}
 \label{eq:cubic-indicial-roots}
\]
These are exactly the two eigenvalues of the canonical modified residue
$R_X$ computed in Section~\ref{sec:atomic-one-step}.  In particular,
Proposition~\ref{prop:residue-discriminant-exponents} identifies their
squared difference with $\delta^\sharp(\alpha_X)=4/9$; the present
calculation independently recovers the same representatives from formal
solutions.
The constant matrix \(C\) adjoins only the algebraic Novikov coefficient
\(r\), and \eqref{eq:cubic-integral-block-gauge} and its inverse use only
integral powers of \(z\).  Hence they preserve each exponent class in
\(\Q/\Z\): multiplication sends \(z^\rho F_X[[z]]\) to itself, and the same
holds for the inverse.  Thus the two classes in
\eqref{eq:cubic-indicial-roots} remain the classes in the original-\(z\)-disc
frame.  By \eqref{eq:universal-exponential}, the repaired framed convention
assigns to a rational exponent \(\rho\) the eigenvalue
\(\operatorname{Exp}_V(\rho)=e^{2\pi i\rho}\).  Therefore these two classes
give
\[
 \operatorname{Exp}_V\left(-\frac16\right)=e^{-\pi i/3},
 \qquad
 \operatorname{Exp}_V\left(-\frac56\right)
 =e^{-5\pi i/3}=e^{\pi i/3}.
 \tag{5.9j}
 \label{eq:cubic-framed-eigenvalues}
\]
This agrees with Cai's rank-two calculation~\cite[p.~6]{Cai}.

Put \(R=\C[r,r^{-1}]\).  For completeness the scalar blocks may equivalently be written as
\[
 z^2\partial_z\widetilde S_\pm
 =\bigl(\pm6r+h_\pm(z)\bigr)\widetilde S_\pm,
 \qquad h_\pm(z)\in z^2R[[z]].
 \tag{5.9k}
 \label{eq:cubic-rank-one-scalar-equations}
\]
Their normalized formal solutions are
\[
 \begin{aligned}
 \widetilde S_+&=e^{-6r/z}u_+(z),
 &u_+(z)&\in1+zR[[z]],\\
 \widetilde S_-&=e^{6r/z}u_-(z),
 &u_-(z)&\in1+zR[[z]].
 \end{aligned}
 \tag{5.9l}
 \label{eq:cubic-rank-one-solutions}
\]
Indeed, after the displayed exponential is removed,
\(\partial_zu_\pm=z^{-2}h_\pm(z)u_\pm\).  The coefficient
\(z^{-2}h_\pm(z)\) lies in \(R[[z]]\), so integration with initial value
one gives the asserted units.  The exponentials and these units use only
integral powers of the original \(z\); neither adjoins a root of \(z\) nor
produces a deck factor.  Thus each block is unramified and has framed
regular-monodromy eigenvalue \(1\), so neither contributes a primitive sixth
root.  This agrees
with Cai's scalar solutions~\cite[pp.~5--6]{Cai}.

It remains to count the rank-two solutions.  Write its ansatz as
\(\widetilde S_3=z^\rho\sum_{n\ge0}a_nz^n\) and
\(\widetilde S_4=z^{\rho+1}\sum_{n\ge0}b_nz^n\), with \(a_0=1\).
At coefficient \(n\), the two equations have the form
\[
 L_{\rho+n}\begin{pmatrix}a_n\\b_n\end{pmatrix}
 =\begin{pmatrix}\alpha_n\\\beta_n\end{pmatrix},
 \qquad
 L_s=
 \begin{pmatrix}
 s+19/18&-2\\
 8/81&s-1/18
 \end{pmatrix},
 \qquad
 \det L_s=\left(s+\frac16\right)\left(s+\frac56\right).
 \tag{5.9m}
 \label{eq:cubic-frobenius-recursion}
\]
Here \((\alpha_n,\beta_n)\) depends only on coefficients of lower index;
for \(n=0\) it is zero, and \eqref{eq:cubic-indicial-polynomial} is exactly
the condition that \(L_\rho\) have a kernel.  The roots in
\eqref{eq:cubic-indicial-roots} are distinct modulo \(\Z\), so
\(L_{\rho+n}\) is invertible for every \(n\ge1\).  For either root,
\(\ker L_\rho\) is one-dimensional and its first coordinate is nonzero;
the normalization \(a_0=1\) fixes \(b_0\).  The recursion therefore gives
one formal solution for each root and no further solution in its exponent
class.  The two resulting solutions are independent, and exhaust the
rank-two block.

Consequently each eigenvalue in \eqref{eq:cubic-framed-eigenvalues} has
algebraic multiplicity one, while
\eqref{eq:cubic-rank-one-solutions} contributes neither primitive sixth root.
Hence
\[
 \nu_6(X)=2.
 \tag{5.9n}
 \label{eq:cubic-packet-internal-count}
\]
This is a calculation in the small even connection, so no odd-cohomology block
enters the count.  Moreover, \(N_1(X)=\Z\ell\), and the numerical Novikov
ring is exactly \(\C[[Q^\ell]]=\C[[q]]\); thus the numerical-Novikov
convention does not alter the calculation.

The shared reduction in Proposition~\ref{prop:cubic-block-data} is proved
internally from Cai's displayed small-even matrix; the argument above uses
only that reduction and the ordinary formal recursion in the original loop
coordinate.

\end{proof}

\begin{corollary}[Formal-germ persistence of the cubic packet]
\label{cor:cubic-formal-germ}
After extending the cubic Novikov coefficient field so that
$r=(3q)^{1/2}$ is a unit, the primitive-sixth multiplicity is equal to $2$
throughout the formal even bulk germ at the small hyperplane point.
\end{corollary}

\begin{proof}
At the closed point the Euler eigenvalues are $6r,-6r,0$, with the zero
cluster a single rank-two Jordan block and the two nonzero clusters of rank
one.  The differences from the zero eigenvalue are units, so
Proposition~\ref{prop:ranktwo-framed-germ} applies to the zero packet and
freezes its two exponent classes $-1/6,-5/6$.  The two simple packets remain
rank one on the formal germ and contribute no primitive sixth root by
Lemma~\ref{lem:simple-euler-block}.  Proposition~\ref{prop:cubic-packet}
sets the value at the closed point equal to $2$.
\end{proof}

\begin{corollary}[Framed count after one product stabilization]
\label{cor:cubic-product-nu}
For every smooth cubic threefold $X$,
\[
 \nu_6(X\times\PP^1)=4.
\]
This statement is unconditional.
\end{corollary}

\begin{proof}
Combine Proposition~\ref{prop:projective-product-nu} with
Proposition~\ref{prop:cubic-packet}.
\end{proof}

\begin{remark}[Conditional second proof of Theorem~\ref{thm:every-cubic}]
Assuming Hypotheses~5.7R and 5.7T,
Theorem~\ref{thm:nu6-birational-invariance} turns the unconditional value
\(\nu_6(X\times\PP^1)=4\) from Corollary~\ref{cor:cubic-product-nu} into
the original framed-monodromy proof: a rational smooth fourfold has
\(\nu_6=0\).  This remains a second, stronger-invariant proof rather than
the logical basis of Theorem~\ref{thm:every-cubic}.
\end{remark}

\subsection{A genus-eight Fano consequence}

Kuznetsov associates to every smooth prime Fano threefold \(V\) of genus
eight a smooth Pfaffian cubic threefold \(X\), together with rank-two vector
bundles \(U\) on \(V\) and \(E^*\) on \(X\), whose projectivizations are
related by a flop \cite[Theorems~2.17--2.18]{KuznetsovV14}.

\begin{corollary}[One stabilization of \(V_{14}\)]
\label{cor:v14-one-step}
For every smooth prime Fano threefold \(V\) of genus eight,
\[
  V\times\PP^1\ \text{is irrational}.
\]
Assuming Hypotheses~5.7R and~5.7T, one moreover has
\[
  \nu_6(V)=2.
  \tag{5.11}
  \label{eq:v14-stabilization}
\]
\end{corollary}

\begin{proof}
The flop gives a birational map between \(\PP_V(U)\) and
\(\PP_X(E^*)\).  A projective bundle is birational to the product of its
base with the corresponding projective space, so
\(V\times\PP^1\) is birational to \(X\times\PP^1\).  The latter is
irrational by Theorem~\ref{thm:every-cubic}, proving the unconditional
assertion.

Assume now Hypotheses~5.7R and~5.7T.  Both projectivizations are smooth
projective
fourfolds.  By Theorem~\ref{thm:nu6-birational-invariance}, Proposition
\ref{prop:framed-operations}, and Proposition~\ref{prop:cubic-packet},
\[
  2\nu_6(V)=\nu_6(\PP_V(U))
  =\nu_6(\PP_X(E^*))=2\nu_6(X)=4.
\]
Hence \(\nu_6(V)=2\).
\end{proof}

The two-hypothesis framed refinement is sharp at this level.  In dimension five
a weak-factorization center may have dimension three and may itself carry
the cubic packet, so the low-dimensional vanishing used in
Theorem~\ref{thm:nu6-birational-invariance} no longer prevents
cancellation.  The ordinary atom obstruction has the parallel limitation
recorded after Theorem~\ref{thm:every-cubic}: from the second stabilization
onward the cubic atom is already allowed as a lower-dimensional
representative.
